\RAsection{RA6}
\subsection*{Resultados de aprendizaje}
Analizar el diseño y la construcción de módulos de programa para la resolución de problemas complejos considerando como un medio eficaz la división de programas en componentes más sencillos que interactúan entre sí.

\subsection*{Contenido según programa}
Funciones en lenguaje C.  Programación modular.  Definición de una función.  Prototipo de una función.  Argumentos. Archivos de cabecera. Variables locales y globales. Clases de Almacenamiento: externas, estáticas, registros y automáticas. Reglas de alcance (scope).  Proposición return. Llamada por valor y por referencia.  Recursividad. Condiciones para implementar funciones recursivas. Ejemplo: factorial de un número.  Recursividad vs. Iteración.

%--------------------------------------------------------------------------------
%--------------------------------------------------------------------------------
\setcounter{subsection}{6}
\subsection*{Funciones $1^{era}$ Parte}

\subsubsection{Ejercicio}
Completar el siguiente programa:

\lstset{inputencoding=utf8/latin1}
\lstinputlisting[style=customc]{code/ejercicio1_fun_completar_guia.c}
{\small
  \lstset{inputencoding=utf8/latin1}
  \lstinputlisting[backgroundcolor = \color{lightgray}]{code/ejercicio1_completar_salida.mn}
}

\subsubsection{Ejercicio}
Completar el siguiente programa:
\lstset{inputencoding=utf8/latin1}
\lstinputlisting[style=customc]{code/ejercicio2_completar_guia.c}
Para obtener la siguiente salida:
{\small
  \lstset{inputencoding=utf8/latin1}
  \lstinputlisting[backgroundcolor = \color{lightgray}]{code/ejercicio2_completar_salida.mn}
}

\subsubsection{Ejercicio}
Completar el siguiente programa:
\lstset{inputencoding=utf8/latin1}
\lstinputlisting[style=customc]{code/ejercicio3_completar_guia.c}
Para obtener la siguiente salida:
{\small
  \lstset{inputencoding=utf8/latin1}
  \lstinputlisting[backgroundcolor = \color{lightgray}]{code/ejercicio3_completar_salida.mn}
}

\subsubsection{Ejercicio}
Completar el siguiente programa:
\lstset{inputencoding=utf8/latin1}
\lstinputlisting[style=customc]{code/ejercicio4_completar_guia.c}
Para obtener la siguiente salida:
{\small
  \lstset{inputencoding=utf8/latin1}
  \lstinputlisting[backgroundcolor = \color{lightgray}]{code/ejercicio4_completar_salida.mn}
}

\subsubsection{Ejercicio}
Completar el siguiente programa:
\lstset{inputencoding=utf8/latin1}
\lstinputlisting[style=customc]{code/ejercicio5_completar_guia.c}
Para obtener la siguiente salida:
{\small
  \lstset{inputencoding=utf8/latin1}
  \lstinputlisting[backgroundcolor = \color{lightgray}]{code/ejercicio5_completar_salida.mn}
}


\subsubsection{Ejercicio}
Completar el siguiente programa:
\lstset{inputencoding=utf8/latin1}
\lstinputlisting[style=customc]{code/ejercicio6_completar_guia.c}
Para obtener la siguiente salida:
{\small
  \lstset{inputencoding=utf8/latin1}
  \lstinputlisting[backgroundcolor = \color{lightgray}]{code/ejercicio6_completar_salida.mn}
}

\subsubsection{Ejercicio}
Completar el siguiente programa:
\lstset{inputencoding=utf8/latin1}
\lstinputlisting[style=customc]{code/ejercicio7_completar_guia.c}
Para obtener la siguiente salida:
{\small
  \lstset{inputencoding=utf8/latin1}
  \lstinputlisting[backgroundcolor = \color{lightgray}]{code/ejercicio7_completar_salida.mn}
}

\subsubsection{Ejercicio}
Crear los prototipos para las diferentes funciones que resuelvan:
\begin{itemize}
  \item Función \textbf{hipotenusa}, toma dos argumentos float de doble precisión, \textbf{lado1} y \textbf{lado2}. Devuelve el resultado en float de doble precisión.
  \item Función \textbf{menor}, toma tres enteros, \textbf{x}, \textbf{y}, \textbf{z}. Devuelve un entero.
  \item Función \textbf{intToFloat}, que toma un argumento entero llamado \textbf{numero}, y devuelve el resultado en float.
\end{itemize}

\subsubsection{Calculadora con operaciones básicas}
Implementar una función para cada operación matemática básica. Cada función debe recibir dos números reales y devolver el resultado correspondiente. La función dividir debe verificar que el divisor no sea cero. Si lo es, retornar 0 e imprimir un mensaje de error desde la función.

Prototipos:

\begin{lstlisting}[style=customc, basicstyle=\ttfamily]
float sumar(float a, float b);
float restar(float a, float b);
float multiplicar(float a, float b);
float dividir(float a, float b);
\end{lstlisting}


\subsubsection{Verificar si un número es múltiplo de otro}
La función debe devolver 1 si a es múltiplo de b, y 0 en caso contrario. Considerar el caso en el que b sea 0 e imprimir un mensaje de error en ese caso.
\begin{lstlisting}[style=customc, basicstyle=\ttfamily]
int esMultiplo(int a, int b);
\end{lstlisting}

\subsubsection{Clasificar un carácter}
La función debe indicar si el carácter recibido es una letra mayúscula, letra minúscula, dígito numérico (0-9) o otro símbolo. Utilizar condicionales y switch cuando sea posible.
\begin{lstlisting}[style=customc, basicstyle=\ttfamily]
void clasificarCaracter(char c);
\end{lstlisting}

\subsubsection{Contador de cifras}
La función debe contar cuántas cifras tiene un número entero positivo. Por ejemplo, si recibe 1234 debe devolver 4. Si recibe 0 debe devolver 1. No debe usar strings ni arreglos.
\begin{lstlisting}[style=customc, basicstyle=\ttfamily]
int contarCifras(int n);
\end{lstlisting}

\subsubsection{Contador de dígitos pares en un número}
La función debe contar cuántos dígitos pares tiene un número entero positivo. No se pueden usar arreglos. Por ejemplo, si el número es 248, debe devolver 3.
\begin{lstlisting}[style=customc, basicstyle=\ttfamily]
int contarPares(int n);
\end{lstlisting}

\subsubsection{Conversor de días a semanas y días}
La función debe calcular cuántas semanas y días hay en una cantidad total de días. Debe imprimir por pantalla los resultados: semanas y dias.
\begin{lstlisting}[style=customc, basicstyle=\ttfamily]
void convertirDias(int totalDias);
\end{lstlisting}

\subsubsection{Determinar si un año es bisiesto}
La función debe devolver 1 si el año es bisiesto y 0 si no lo es. Un año es bisiesto si es múltiplo de 4 y (no es múltiplo de 100 o sí es múltiplo de 400).
\begin{lstlisting}[style=customc, basicstyle=\ttfamily]
int esBisiesto(int anio);
\end{lstlisting}


\subsubsection{Cajero Automático Simplificado}
Simular el retiro de dinero de un cajero automático.

Primero, el usuario debe ingresar un PIN. Si no es el correcto (por ejemplo, 1234), el sistema debe mostrar un mensaje de error y no continuar.

Si el PIN es correcto, el usuario ingresa un monto a retirar. El sistema debe calcular cuántos billetes de una denominación específica (por ejemplo, billetes de 100) se entregarán.

Mostrar cuántos billetes se entregan y cuál es el saldo restante (que no se puede entregar porque no alcanza para un billete entero).
\begin{lstlisting}[style=customc, basicstyle=\ttfamily]
int verificarPin(int pinIngresado);
int calcularBilletes(int monto, int denominacion);
\end{lstlisting}


\subsubsection{Validar Horario}
Crear una función que reciba tres enteros: horas, minutos y segundos. Verificar que:
Las horas estén entre 0 y 23.
Los minutos y segundos estén entre 0 y 59.
Si los valores son válidos, devolver 1. Si no, devolver 0 y mostrar un mensaje de error indicando cuál campo es incorrecto.
\begin{lstlisting}[style=customc, basicstyle=\ttfamily]
int validarHora(int horas, int minutos, int segundos); 
\end{lstlisting}

\subsubsection{Convertir Minutos Totales a Horas y Minutos}
Dada una cantidad total de minutos, calcular cuántas horas completas hay y cuántos minutos sobran.
El valor de minutos debe ser mayor o igual a 0. Si es negativo, mostrar un mensaje de error y pedir de nuevo el valor.
Los resultados deben devolverse a través de parámetros por referencia.
\begin{lstlisting}[style=customc, basicstyle=\ttfamily]
void convertirTiempo(int minutosTotales, int* horas, int* minutos);
\end{lstlisting}

\subsubsection{Calculadora de Plazos}
Desarrollar un programa para calcular la cantidad de cuotas que tiene que pagar un cliente al comprar un producto en diferentes plazos. El programa debe:
\begin{itemize}
  \item Pedir al usuario el precio del producto (no puede ser negativo).
  \item Preguntar cuántas cuotas quiere realizar (no puede ser menor a 2 ni mayor a 12).
  \item Pedir la tasa de interés anual (no puede ser negativa).
  \item Calcular el monto de la primera cuota.
\end{itemize}

\begin{lstlisting}[style=customc, basicstyle=\ttfamily]
\end{lstlisting}

\subsubsection{}
\begin{lstlisting}[style=customc, basicstyle=\ttfamily]
\end{lstlisting}

%--------------------------------------------------------------------------------
%--------------------------------------------------------------------------------
\subsection*{Funciones $2^{da}$ Parte}

\subsubsection{Ejercicio}
Implementar la función \texttt{fibonacci} en forma recursiva.

\lstset{inputencoding=utf8/latin1}
\lstinputlisting[style=customc]{code/ejercicio1_fun2_completar_guia.c}

\subsubsection{Ejercicio}
El rompecabezas de la Torre de Hanoi fue inventado por el matemático francés Edouard Lucas en 1883. Se inspiró en una leyenda acerca de un templo hindú donde el rompecabezas fue presentado a los jóvenes sacerdotes. Al principio de los tiempos, a los sacerdotes se les dieron tres postes y una pila de 64 discos de oro, cada disco un poco más pequeño que el de debajo. 

Su misión era transferir los 64 discos de uno de los tres postes a otro, con dos limitaciones importantes. Sólo podían mover un disco a la vez, y nunca podían colocar un disco más grande encima de uno más pequeño. Los sacerdotes trabajaban muy eficientemente, día y noche, moviendo un disco cada segundo. Cuando terminaran su trabajo, dice la leyenda, el templo se desmenuzaría en polvo y el mundo se desvanecería.

En la Figura  Se puede ver la representación gráfica del juego.

% \begin{figure}[h!]
% \centering
% \includegraphics[width=\textwidth]{img/hanoi.eps}
% \caption{Torres de Hanoi\\ Fuente: Deitel\&Deitel}
% \label{fig:hanoi}
% \end{figure}
Supongamos que los sacerdotes intentan mover los discos de la clavija 1 a la clavija 3. Deseamos desarrollar un algoritmo que imprima la secuencia precisa de transferencias de clavija de disco a disco.
Si tuviéramos que abordar este problema con métodos convencionales, nos encontraríamos rápidamente anudados en la gestión de los discos. En cambio, si atacamos el problema con la recursión en mente, inmediatamente se vuelve manejable. Mover n discos se puede ver en términos de mover solo n - 1 discos (y, por lo tanto, la recursividad) de la siguiente manera:
\begin{enumerate}
  \item Mueva n - 1 discos de la clavija 1 a la clavija 2, utilizando la clavija 3 como área de retención temporal.
  \item Mueva el último disco (el más grande) de la clavija 1 a la clavija 3.
  \item Mueva los discos n - 1 de la clavija 2 a la clavija 3, utilizando la clavija 1 como área de retención temporal.
\end{enumerate}
El proceso finaliza cuando la última tarea implica mover n = 1 disco, es decir, el caso base. Esto se logra moviendo trivialmente el disco sin la necesidad de un área de retención temporal.
Escribe un programa para resolver el problema de Towers of Hanoi. Use una función recursiva con cuatro parámetros:
\begin{enumerate}
  \item El número de discos a mover
  \item La clavija en la que se enroscan estos discos inicialmente
  \item La clavija a la que se debe mover esta pila de discos
  \item La clavija que se utilizará como área de retención temporal.
\end{enumerate}
Su programa debe imprimir las instrucciones precisas que necesitará para mover los discos desde la clavija de inicio a la clavija de destino. Por ejemplo, para mover una pila de tres discos de la clavija 1 a la clavija 3, su programa debe imprimir la siguiente serie de movimientos:

{\scriptsize
  \lstset{inputencoding=utf8/latin1}
  \lstinputlisting[backgroundcolor = \color{lightgray}]{code/hanoi_salida.mn}
}

\subsection*{\underline{Ejercicios de cierre}}
\subsubsection{Juego del ahorcado}
Escriba un programa que implemente el juego del ahorcado.
El programa debe solicitar al usuario que adivine una palabra desconocida ingresando letras una por una.
Si la letra está en la palabra, el programa debe mostrar su ubicación en la palabra.
Si la letra no está en la palabra, el programa debe mostrar un mensaje indicando que la letra no es correcta y dibujar una parte del ahorcado.
El juego continúa hasta que el usuario adivine la palabra completa o se complete el dibujo del ahorcado.

El siguiente código puede servir de ayuda.

\lstset{inputencoding=utf8/latin1}
\lstinputlisting[style=customc]{code/ahorcado.c}
